% 第5章：元启发式与优化方法
\chapter{元启发式与优化方法}

\begin{levelbox}
元启发式方法在组合优化和任务调度中有广泛应用。遗传算法和模拟退火（5.1--5.2节）为本科生内容，群体智能算法（5.3节）为研究生内容，新型元启发式（5.4节）为自学拓展。
\end{levelbox}

% =============================================
\section{进化算法 \ug}

\subsection{遗传算法（GA）}

遗传算法模拟自然进化过程，通过选择、交叉和变异操作搜索解空间。

\begin{algbox}[title={\textbf{遗传算法基本框架}}]
\begin{enumerate}
  \item \textbf{编码：}将任务序列/资源分配方案编码为染色体（如排列编码、二进制编码）；
  \item \textbf{初始化：}随机生成初始种群；
  \item \textbf{适应度评估：}计算每个个体的目标函数值；
  \item \textbf{选择：}按适应度选择父代（轮盘赌、锦标赛选择等）；
  \item \textbf{交叉：}两个父代交换部分基因产生后代（顺序交叉OX、部分匹配交叉PMX等）；
  \item \textbf{变异：}以小概率改变基因值（交换变异、插入变异等）；
  \item \textbf{迭代：}重复直到满足终止条件。
\end{enumerate}
\end{algbox}

\textbf{在规划中的应用：}GA特别适合任务调度和资源分配问题。例如，Job-Shop调度问题中，将工序排列编码为染色体，使用最小完工时间（Makespan）作为适应度函数。

\begin{appbox}[title={\textbf{应用：车辆路径规划（VRP）}}]
在物流配送中，需要为一组车辆规划服务路线以最小化总行驶距离。GA可以有效求解大规模VRP问题：将每条路线编码为客户访问序列，使用特殊的交叉算子（如NWOX）保持路线的合法性。
\end{appbox}

\subsection{差分进化（DE）\grad}

差分进化使用种群中个体间的差异向量来驱动搜索：
$$\mathbf{v}_i = \mathbf{x}_{r_1} + F \cdot (\mathbf{x}_{r_2} - \mathbf{x}_{r_3})$$

DE在连续优化问题中表现优异，参数少且鲁棒性强。在连续域任务规划（如机器人轨迹优化）中有良好的应用。

% =============================================
\section{模拟退火与禁忌搜索 \ug}

\subsection{模拟退火（SA）}

\begin{algbox}[title={\textbf{模拟退火核心思想}}]
类比金属退火过程：
\begin{enumerate}
  \item 从初始解和初始高温$T_0$开始；
  \item 在当前解的邻域中随机选取新解$s'$；
  \item 若$f(s') < f(s)$（改善），接受$s'$；
  \item 若$f(s') \geq f(s)$（恶化），以概率$\exp(-\Delta f / T)$接受$s'$；
  \item 按冷却计划降低温度$T$；
  \item 重复直到温度足够低。
\end{enumerate}
\textbf{关键：}高温时容易接受较差解（探索），低温时趋向于只接受改善（开发）。
\end{algbox}

\subsection{禁忌搜索（TS）\grad}

禁忌搜索使用\textbf{禁忌表}（Tabu List）记录最近访问过的解（或移动），禁止短期内回访以跳出局部最优。

核心组件：
\begin{itemize}
  \item \textbf{禁忌表：}存储最近$k$步的移动，禁止这些移动的逆操作；
  \item \textbf{特赦准则（Aspiration Criteria）：}若某禁忌移动产生的解优于历史最优，则解除禁忌；
  \item \textbf{邻域结构：}定义从当前解到邻近解的转换操作。
\end{itemize}

% =============================================
\section{群体智能算法 \grad}

\subsection{粒子群优化（PSO）}

粒子群优化模拟鸟群觅食行为。每个粒子代表一个候选解，根据自身历史最优和群体历史最优更新位置：
$$\mathbf{v}_i^{t+1} = w\mathbf{v}_i^t + c_1 r_1 (\mathbf{p}_i - \mathbf{x}_i^t) + c_2 r_2 (\mathbf{g} - \mathbf{x}_i^t)$$
$$\mathbf{x}_i^{t+1} = \mathbf{x}_i^t + \mathbf{v}_i^{t+1}$$

\textbf{离散PSO：}标准PSO处理连续变量。用于组合优化时需要离散化处理，如将位置映射为排列（排列PSO）或使用集合操作定义速度更新。

\subsection{蚁群优化（ACO）}

蚁群优化模拟蚂蚁通过信息素寻找最短路径的行为：

\begin{enumerate}
  \item 每只蚂蚁根据信息素浓度和启发信息概率性地构建解；
  \item 构建完成后，根据解的质量更新信息素——好的解路径上信息素增加；
  \item 同时全局信息素以一定速率蒸发，避免过早收敛。
\end{enumerate}

ACO在TSP和路径规划中表现出色，其最优蚁群系统（MMAS）和蚁群系统（ACS）是两个重要变体。

\begin{appbox}[title={\textbf{应用：无人机航路规划}}]
ACO非常适合无人机航路规划：将航路点作为图的节点，信息素表示路径的优劣，启发信息包含距离、威胁和能耗等。蚂蚁从起点出发，概率性选择下一个航路点，最终找到安全高效的飞行路线。
\end{appbox}

% =============================================
\section{新型元启发式算法 \self}

近年来涌现了大量受自然现象启发的新型元启发式算法：

\subsection{灰狼优化（GWO）}

模拟灰狼群的社会等级和捕猎行为。狼群分为$\alpha$（最优）、$\beta$（次优）、$\delta$（第三）和$\omega$（其余），搜索过程模拟包围和攻击猎物。

\subsection{鲸鱼优化算法（WOA）}

模拟座头鲸的气泡网捕食策略，包含三种行为：包围猎物、螺旋游动和随机搜索。

\subsection{哈里斯鹰优化（HHO）}

模拟哈里斯鹰的协同捕猎策略，包含探索阶段和开发阶段（突袭行为）。

\begin{remark}
新型元启发式算法数量众多（100+种），但很多在本质上与已有算法相似。在选择使用时，应关注：(1) 是否有严格的理论分析；(2) 是否在标准基准上与经典方法进行了公平比较；(3) 是否解决了实际问题而非仅为新颖性。
\end{remark}

% =============================================
\section{元启发式在任务规划中的应用框架}

\subsection{编码策略}

将规划问题转化为优化问题的关键是设计合适的编码：
\begin{itemize}
  \item \textbf{排列编码：}适用于调度问题（任务执行顺序）
  \item \textbf{二进制编码：}适用于选择问题（是否执行某动作）
  \item \textbf{实数编码：}适用于连续参数优化（轨迹控制点）
  \item \textbf{混合编码：}同时包含离散和连续决策变量
\end{itemize}

\subsection{约束处理}

规划问题通常包含大量约束。常用的约束处理技术：
\begin{itemize}
  \item 罚函数法：将约束违反度加入目标函数
  \item 修复算子：对不可行解进行修复
  \item 解码器法：设计解码器保证总是生成可行解
\end{itemize}

% =============================================
\section{本章小结与习题}

\subsection*{本章小结}
\begin{itemize}
  \item 元启发式方法适合处理NP-hard的组合优化和调度问题，不依赖问题的特殊结构。
  \item GA和SA是最经典的元启发式，易于理解和实现。
  \item 群体智能算法（PSO、ACO）通过群体协作搜索，在路径规划等问题上有独特优势。
  \item 选择元启发式时应注重实效性而非新颖性。
\end{itemize}

\subsection*{思考与练习}
\begin{enumerate}
  \item 用遗传算法求解一个5任务3机器的Job-Shop调度问题。
  \item 比较模拟退火和禁忌搜索在TSP问题上的求解质量和收敛速度。
  \item （研究生）用ACO算法解决一个带时间窗的车辆路径问题（VRPTW）。
  \item （研究生）分析PSO和GWO在数学本质上的异同。
\end{enumerate}
